\chapter*{What the Prestige Still Owes}
\addcontentsline{toc}{chapter}{What the Prestige Still Owes}

\noindent\textbf{[This section records what remains unwritten in the
third part of \emph{Finding Mind}.]}

\bigskip

The Prestige --- in the tradition of the magic trick whose structure
organizes this book --- is where the impossible is brought back. The
coin vanishes in the Pledge; it performs the extraordinary in the
Turn; in the Prestige, it reappears, and we are invited to understand
both what happened and why we could not see it happening.

Three of its chapters now exist.
Chapter~\ref{ch:kant} reads the whole construction as a
transcendental philosophy made internal.
Chapter~\ref{ch:trick} shows the method: the restriction of scope to
computation, which yielded an ontology and a grounded hypothesis
language for free, and which is the sleight of hand underneath every
result in the Turn.
Chapter~\ref{ch:notation} is the resulting price list --- what a
symbol system is, which of Goodman's conditions the machinery
supplies and which it does not, what a budget buys, and how far the
calculus-making functor $\RHO[-]$ can be pushed before it consumes
what it was meant to produce.

\bigskip

What is still owed:

\begin{enumerate}
  \item \textbf{The View from Inside} --- What a conscious agent at
        level $\omega$ of the hypercomputational fibration can and
        cannot know about its own substrate; the first-person /
        third-person interface as a theorem rather than a mystery.
        Chapter~\ref{ch:trick} supplies part of the negative half of
        this --- what such an agent cannot know is bounded below by
        what its population has not collectively grounded --- but the
        positive account is unwritten.
  \item \textbf{The Recapitulation} --- Origin of mind as origin of
        life one octave up: the replication of fixed points, the
        emergence of self-models, and the evolutionary pressure
        toward accurate internal simulation. The existence proofs of
        Section~\ref{sec:soluble} belong to this chapter as much as
        to the one they appear in, and the assembly-index arguments
        of Part~\ref{part:origins} are the other half of it.
  \item \textbf{The Metric} --- The problem
        Section~\ref{sec:kant-coda} names and this book does not
        solve: a notion of distance under which computational
        behaviors that are nearly the same come out near each other,
        without which the space read off a term structure stays
        merely combinatorial.
        Observation~\ref{obs:metric} is a warning about how not to do
        it, not a construction.
  \item \textbf{Finding Mind} --- The synthesis: what it means to
        find mind in the structure of computation, and what that
        finding demands of us --- scientifically, philosophically,
        and ethically. Section~\ref{sec:blind-spot-returned} states
        the constraint this chapter has to work under: a mind that
        acts in the world we act in either shares our correspondence,
        and is therefore a member of our population in a
        constitutive rather than a decorative sense, or builds its
        own with a population of its own. There is no third
        arrangement. The ethics is downstream of that, and it has not
        been written.
\end{enumerate}

\bigskip
\noindent\hrule
\bigskip

\noindent\textit{The remainder of the Prestige is forthcoming.
Readers wishing to follow its development are invited to the
author's Substack at \texttt{https://f1r3fly.substack.com}.}
